\documentclass[11pt,a4paper]{article}

\usepackage[utf8]{inputenc}
\usepackage[T1]{fontenc}
\usepackage[margin=25mm]{geometry}
\usepackage{amsmath,amssymb}
\usepackage{booktabs}
\usepackage{longtable}
\usepackage{xcolor}
\usepackage{listings}
\usepackage{enumitem}
\usepackage[colorlinks=true,linkcolor=NavyAccent,urlcolor=NavyAccent,
            citecolor=NavyAccent]{hyperref}

\definecolor{NavyAccent}{HTML}{26405E}
\definecolor{CodeBg}{HTML}{F2F5F8}
\definecolor{Rule}{HTML}{C8CCD2}
\definecolor{Alert}{HTML}{A8322D}
\definecolor{Good}{HTML}{1F6F43}
\definecolor{Muted}{HTML}{5B6470}

\lstset{
  basicstyle=\ttfamily\footnotesize,
  backgroundcolor=\color{CodeBg},
  frame=single,
  rulecolor=\color{Rule},
  breaklines=true,
  columns=fullflexible,
  keepspaces=true,
  showstringspaces=false,
  xleftmargin=0pt,
  framexleftmargin=4pt,
  framexrightmargin=4pt,
  aboveskip=8pt,
  belowskip=8pt,
}

% Long ttfamily identifiers cannot hyphenate; let TeX stretch the line instead of
% overflowing the margin.
\emergencystretch=3em

\newcommand{\jl}[1]{\texttt{\small #1}}
\newcommand{\file}[1]{\texttt{\small #1}}
\newcommand{\finding}[1]{\textcolor{Alert}{\textbf{#1}}}
\newcommand{\ok}[1]{\textcolor{Good}{#1}}
\newcommand{\gnorm}{\lVert g\rVert}
\newcommand{\gk}{\lVert g_k\rVert}

\title{\vspace{-14mm}\textcolor{NavyAccent}{\textbf{What experiments 5, 7 and 9 do}}}
\author{\normalsize A reading of \file{exp5\_inactivity.jl},
        \file{exp7\_convergence\_rate.jl} and \file{exp9\_second\_order.jl}}
\date{\normalsize Read against branch \file{checking-w-Claude-Code}, 2026-08-20}

\begin{document}
\maketitle
\vspace{-6mm}

\noindent\textcolor{Muted}{\small All three files are read as source and
cross-checked against their most recent archived run (2026-08-19, 16:28). Where
the code and the archive disagree, the disagreement is reported rather than
reconciled.}

\vspace{2mm}
\hrule
\vspace{4mm}

\section*{Summary in one paragraph each}

\textbf{Experiment 5} measures whether the trust-region constraint ever
\emph{stops} binding. The claim is that the mechanisms split cleanly into two
groups by this property, and -- the part that matters -- that the split is
\emph{invisible} to a first-order convergence test, because both groups drive
$\gk \to 0$ perfectly well. The observable is the fraction of late iterations with
$\lVert s_k\rVert = \Delta_k$. Two configurations are included with parameters
deliberately set below any plausible threshold, as a positive control: they should
stay pinned at 1 while every other rule falls to 0.

\smallskip
\textbf{Experiment 7} measures the \emph{observed local convergence order} of each
radius mechanism and tests a negative claim: that the order does \emph{not} depend
on the mechanism. Once the trust-region constraint stops binding, the step is the
unconstrained model minimiser, so the rate is a property of the model rather than
of the radius rule. What the mechanism controls is \emph{when} that happens. The
script therefore fits a rate only on the inactive tail of each run, and reports the
fitted order beside $k_\star$, the first permanently inactive iteration.

\smallskip
\textbf{Experiment 9} tests what happens when a radius rule anchored to $\gnorm$
meets a saddle point. At a saddle $\gnorm = 0$ exactly as at a minimiser, so a
$\gnorm$-anchored rule sets the radius it would use at a solution and the run
stops, reporting success. The script demonstrates the collapse, then shows that
repairing it needs \emph{three} ingredients simultaneously: the criticality measure
$\tau = \max\{\gnorm, -\lambda_{\min}\}$, a subsolver that can move along a
negative-curvature direction, and a model capable of reporting negative curvature
at all.

\section{Experiment 5: direct measurement of inactivity}

\subsection{The claim under test}

Part~II divides the mechanisms by whether the trust-region constraint eventually
stops binding. Experiment 5 measures that division directly, and its real subject
is the second half of the claim:

\begin{quote}
\emph{Both groups drive $\gk \to 0$ monotonically. A first-order convergence test
cannot tell them apart.}
\end{quote}

This is why the experiment exists as a separate script rather than as a column in
the comparison table. A method whose constraint binds for ever still converges, and
still reports \jl{:first\_order}. What it loses is the local rate: the step is never
the unconstrained model minimiser, so the superlinear phase never arrives. Nothing
in $\gk$, in $\rho_k$, or in the exit status reveals this. Only the activity flag
does.

\subsection{The observable, and the floor placed under it}

The observable is \jl{:active\_trajectory}, the per-iteration Boolean
$\lVert s_k\rVert = \Delta_k$ that the solver records when \jl{trace = true}. It is
reduced to a single number per run by \jl{active\_fraction}, the share of the last
10\% of iterations on which the constraint bound.

That reduction is fragile on short runs, and the file guards it:

\begin{lstlisting}
const MIN_ITERS_FOR_TAIL = 30

tail_active(r::RunRecord; frac::Float64 = 0.1) =
    length(r.active_traj) < MIN_ITERS_FOR_TAIL ? NaN :
        active_fraction(RecordView(r); tail = frac)
\end{lstlisting}

The reasoning is in the docstring: on a ten-iteration run the last 10\% is a single
iteration, so \jl{tail\_active} can only be $0$ or $1$ and a table of such values
measures rounding rather than asymptotics. Runs below the floor return \jl{NaN},
are excluded from the mean, and are counted separately in a \jl{measured} column.
Section~\ref{sec:exp5-selection} argues that this guard, while correct, has a
consequence the summary table does not currently account for.

\subsection{The rule list, and two deliberate controls}

Experiment 5 does not use the shared \jl{RULES}. It defines its own six:

\begin{longtable}{@{}p{34mm}p{112mm}@{}}
\toprule
\textbf{configuration} & \textbf{why it is in the list} \\
\midrule
\endhead
\jl{RDelta} & Classical, criticality-blind. Expected to go inactive. \\[2pt]
\jl{RStep} & Step-driven. Expected to go inactive. \\[2pt]
\jl{RDFO($\zeta$=1)} & Criticality-anchored at a sane parameter. \\[2pt]
\jl{RDFO($\zeta$=0.01)} & \textbf{Control.} $\zeta$ far below any plausible
$\bar\kappa$. \\[2pt]
\jl{RGrad} & Gradient-scaled, uncapped. \\[2pt]
\jl{RGradCapped($\mu = \mu_{\max} = 0.05$)} & \textbf{Control.} The cap pins the
multiplier below threshold, so the region can never grow enough to release. \\
\bottomrule
\end{longtable}

The two controls are the design's strongest feature. The closing note says why
explicitly: their parameters sit below any plausible
$\bar\kappa$ ($4/\lambda^*_{\min}$ under \jl{:eigenvalue}, $8/\lambda^*_{\min}$
under \jl{:neighbourhood}), so they \emph{should} stay pinned near 1 while every
other rule falls to 0. An experiment that only shows the expected group succeeding
proves much less than one carrying a case engineered to fail in a known way. If a
control came out inactive, the measurement would be wrong.

\subsection{What the script does}

\begin{enumerate}[leftmargin=7mm,itemsep=3pt]
  \item Opens an archive tagged \jl{inactivity} and records the six rules, the
        shared \jl{SOLVER\_PARAMS} and the problem selection.
  \item Runs the (problem $\times$ rule) grid through \jl{run\_experiment} with
        \jl{trace = true} and, crucially, with\\
        \jl{solver\_kwargs = (hessian\_norm = true,)}.
  \item Reduces each run to a tail-active fraction and aggregates per rule.
  \item Emits the three radius series \emph{per run}, without averaging.
  \item Draws a bar chart of mean tail activity and a cumulative-activity plot.
\end{enumerate}

The \jl{hessian\_norm = true} in step 2 is not incidental. It makes the solver
record $\lVert H_k\rVert$ at every iterate, which is what lets \jl{radius\_sums}
form
\[
  \sum_k \frac{\Delta_k^2}{M_k},
  \qquad M_k = L + \max_{i \le k}\lVert H_i\rVert ,
\]
and the in-line comment states the reason: \emph{that}, not $\sum_k\Delta_k^2$, is
what Part~II proves finite for the criticality-anchored rules. The two agree only
while $\{\lVert H_k\rVert\}$ is bounded, and that boundedness is precisely the
hypothesis the $M_k$ form was introduced to drop. Experiment 5 is currently the
only script in the suite that populates this column.

\subsection{Two refusals to average}

Both are deliberate and both are worth noting, because each removes a way the table
could mislead.

\textbf{No \jl{solved} filter.} The comment is blunt: a configuration trapped below
its threshold is the one that exhausts the budget, so filtering on success would
remove exactly the evidence the experiment exists to collect.

\textbf{No averaging of the radius series across problems.} The second table prints
one row per run. Averaging $\sum_k\Delta_k$ across problems would mix a divergent
series with convergent ones and produce a number with no reading. The reader is
told instead to read the shape across rules \emph{on the same problem}.

\subsection{Outputs}

\begin{longtable}{@{}p{52mm}p{94mm}@{}}
\toprule
\textbf{artefact} & \textbf{content} \\
\midrule
\endhead
\file{exp5\_inactivity\_summary.txt} &
Per rule: runs, solved, \jl{measured}, mean tail activity, \jl{never inactive}.
\textbf{The table the claim lives in.} \\[3pt]
\file{exp5\_radius\_series.txt} &
Per run: iterations, status, $\sum\Delta_k$, $\sum\Delta_k^2$,
$\sum\Delta_k^2/M_k$. \\[3pt]
\file{exp5\_tail\_active.pdf} &
Bar chart of mean tail activity with a dashed line at $0.5$. \\[3pt]
\file{exp5\_active\_}\linebreak[0]\file{trajectory\_AKIVA.pdf} &
Cumulative active fraction along the run. Named for the first problem in the set,
and drawn for that one only. \\
\bottomrule
\end{longtable}

\subsection{What the archived run shows}

From \file{exp\_2026-08-19\_16-28-10\_inactivity}:

\begin{lstlisting}
rule                 runs   solved   measured  mean tail act never inactive
RDelta                  3        3          1          0.000              0
RStep                   3        3          1          0.800              0
RDFO(z=1)               3        3          1          0.000              0
RDFO(z=0.01)            3        3          2          0.600              1
RGrad                   3        3          1          0.400              0
RGradCapped             3        3          3          1.000              3
\end{lstlisting}

\ok{The headline claim is confirmed, and cleanly.} Every configuration solves
every problem -- the \jl{solved} column is $3/3$ throughout -- while
\jl{RGradCapped} never goes inactive on any of them. A first-order test sees six
identical successes; the activity column sees two distinct behaviours. That is the
experiment's whole point, and it lands.

The controls behave as designed. \jl{RGradCapped($\mu = 0.05$)} sits at mean tail
activity $1.000$ with $3/3$ never inactive. \jl{RDFO($\zeta = 0.01$)} is partly
trapped at $0.600$ with $1/3$ never inactive. The cost of being trapped is visible
in the second table: on BEALE, \jl{RDFO($\zeta$=0.01)} takes \textbf{2835}
iterations against \jl{RDFO($\zeta$=1)}'s \textbf{7}, and \jl{RGradCapped} takes
566. Same problem, same tolerance, same first-order status.

The $\sum\Delta_k^2/M_k$ column is populated and spans eight orders of magnitude,
from $6.6\times10^{-5}$ (\jl{RStep} on AKIVA) to $7.4\times10^{3}$ (\jl{RGrad} on
BOXBODLS), which is the shape the section is meant to display.

\subsection{The selection effect in the \jl{measured} column}
\label{sec:exp5-selection}

\finding{Four of the six rules have exactly one measured run.} The mean tail
activity for \jl{RDelta}, \jl{RStep}, \jl{RDFO($\zeta$=1)} and \jl{RGrad} is
computed from a single number, because their other two runs finished in fewer than
\jl{MIN\_ITERS\_FOR\_TAIL = 30} iterations and were excluded.

The exclusion is individually correct -- a one-iteration tail carries no
information -- but the runs it removes are not a random sample. Run length is
\emph{anti-correlated} with the property being measured: a rule that goes inactive
converges quickly and falls below the floor, while a trapped rule grinds on and is
always measured. So the untrapped group is assessed on its slowest, least
representative runs, and the trapped group on all of its.

Two consequences, and they pull in opposite directions.

\begin{itemize}[leftmargin=6mm,itemsep=3pt]
  \item \textbf{The bias is conservative for the headline claim.} A slow untrapped
        run is the one most likely to still be binding late, so measuring only
        those should \emph{inflate} the untrapped group's tail activity. Observing
        \jl{RDelta} and \jl{RDFO($\zeta$=1)} at $0.000$ anyway is therefore stronger
        evidence than it looks, not weaker.
  \item \textbf{It is not conservative for the intermediate values.}
        \jl{RStep} at $0.800$ and \jl{RGrad} at $0.400$ are each one run.
        \jl{active\_fraction}'s own docstring notes that Part~II predicts no
        intermediate value in the limit, so an intermediate value is a finding --
        but on a single run of a few dozen iterations the likelier reading is
        simply that the limit has not been reached. As printed, the table does not
        distinguish those two readings.
\end{itemize}

The \jl{measured} column is what makes the effect visible at all, so the
instrumentation is doing its job. What is missing is that the mean is presented
beside it without qualification. Reporting the count alongside each mean in the
manuscript, or raising the problem set until the floor stops binding, would settle
it. As with experiment 7, the archived run used only three two-variable problems
(AKIVA, BEALE, BOXBODLS), which is the root cause here too: on a larger problem set
most runs would clear 30 iterations honestly.

\section{Experiment 7: observed local convergence order}

\subsection{The claim under test}

Part~II's inactivity theorem says that beyond some iteration $k_\star$ the
constraint $\lVert s_k\rVert \le \Delta_k$ never binds again. Beyond $k_\star$ the
step solves the unconstrained model problem, so the iteration is the underlying
(quasi-)Newton iteration and its rate is whatever that iteration gives. The radius
rule has dropped out.

The experiment turns this into two measurable columns:

\begin{itemize}[leftmargin=6mm,itemsep=2pt]
  \item \textbf{median order} should be roughly \emph{equal} across mechanisms;
  \item \textbf{median $k_\star$} should \emph{differ} substantially, because that
        is the quantity the mechanism decides.
\end{itemize}

A single table showing both is the whole experiment. Note that the claim is a
negative one, which is why the design matters more than usual: an artefact of
estimation would show up as a spurious difference in the first column.

\subsection{Why the estimate must be conditioned on inactivity}

This is the methodological core of the file, and the header says so explicitly.
The order is estimated as the least-squares slope of $\log\lVert g_{k+1}\rVert$ on
$\log\gk$: slope~1 is linear convergence, slope~2 quadratic.

If the fit includes iterations where the constraint was \emph{active}, it mixes a
linear phase (radius-limited steps) with the asymptotic phase (full Newton steps)
and returns a slope between the two. That intermediate number looks like a
convergence order and is not one. It is an artefact of pooling two regimes.

So the script trims every trajectory to its inactive tail before fitting. The
helper is small and worth reading in full:

\begin{lstlisting}
function inactive_tail(r::RunRecord)
    (isempty(r.active_traj) || isempty(r.grad_traj)) && return Float64[]
    k = length(r.active_traj)
    while k >= 1 && !r.active_traj[k]
        k -= 1
    end
    start = k + 1
    start > length(r.active_traj) && return Float64[]
    lo = min(start, length(r.grad_traj))
    return r.grad_traj[lo:end]
end
\end{lstlisting}

It walks backwards from the end of the activity flags over the run of trailing
\jl{false} entries and returns the gradient trajectory from the first of them. If
the constraint was still binding on the final iteration, the loop exits at
\jl{k == length(...)} and the function returns an empty vector, which the caller
treats as the finding for that configuration rather than as missing data.

\subsection{What the script does, in order}

\begin{enumerate}[leftmargin=7mm,itemsep=3pt]
  \item Opens an archive tagged \jl{convergence\_rate} and writes the config,
        recording the rule list, solver parameters and problem selection.
  \item Builds the problem set from \jl{default\_problems()} (CUTEst when
        available, the analytic set otherwise) and the configuration list from
        \jl{rule\_configs()}, which pairs every rule in \jl{RULES} with the default
        model (\jl{ExactHessian}) and default subsolver (\jl{SteihaugCG}).
  \item \textbf{Overrides the tolerance to} $10^{-10}$. This is deliberate and is
        the one parameter the experiment does not inherit: the asymptotic regime is
        precisely what is being measured, so the run must be allowed deep into it.
        The other four parameters are taken from \jl{SOLVER\_PARAMS} unchanged.
  \item Runs everything through \jl{run\_experiment} with \jl{trace = true}, so each
        result is a \jl{RunRecord} carrying the full trajectories.
  \item For each record, computes three quantities and prints a row.
  \item Aggregates per rule into a second table, and draws two figures.
\end{enumerate}

\subsection{The three quantities per run}

\begin{longtable}{@{}p{34mm}p{112mm}@{}}
\toprule
\textbf{quantity} & \textbf{meaning} \\
\midrule
\endhead
\jl{ki} $=k_\star$ &
\jl{inactivity\_index(RecordView(r))}: the first iteration after which the
constraint \emph{never} binds again. \jl{nothing} if it bound to the end. This is
the quantity the mechanism controls. \\[3pt]
\jl{fi} &
\jl{findfirst(!, r.active\_traj)}: the first inactive iteration \emph{of any kind}.
Printed as \jl{fi - 1} to count from zero. It differs from $k_\star$ whenever the
region releases and then binds again, and printing both makes that visible. \\[3pt]
\jl{ord} &
\jl{estimate\_convergence\_order(tail)}, the least-squares slope described above,
computed only when the tail has at least \jl{MIN\_TAIL = 6} points. \\
\bottomrule
\end{longtable}

Each run is then classified into exactly one of three buckets, and this
classification is the part that keeps the summary honest:

\begin{itemize}[leftmargin=6mm,itemsep=2pt]
  \item \jl{never} -- the constraint bound to the end, so there is no asymptotic
        phase to fit;
  \item \jl{short} -- fewer than six inactive iterations, too few for a regression
        that means anything;
  \item \emph{fitted} -- contributes its order and its $k_\star$ to the medians.
\end{itemize}

Only the third bucket enters the medians, and the counts of the other two are
printed beside them. The file states the reason directly: a median taken over the
runs that reached inactivity early is biased towards the mechanisms that get there
early. Without those two counts the summary would flatter exactly the rules the
experiment is trying to distinguish.

There is also a deliberate omission. The loop applies \textbf{no} \jl{solved}
filter. A run that never goes inactive is not dropped; it is counted in
\jl{n\_never} and reported. As the comment puts it, that is the result for the
configuration rather than a gap in the table.

\subsection{Outputs}

\begin{longtable}{@{}p{52mm}p{94mm}@{}}
\toprule
\textbf{artefact} & \textbf{content} \\
\midrule
\endhead
\file{exp7\_conv\_order\_summary.txt} &
One row per (problem, rule): status, iterations, $k_\star$, first inactive, tail
length, fitted order. \\[3pt]
\file{exp7\_order\_by\_rule.txt} &
One row per rule: how many runs were fitted, short and never; median order; median
$k_\star$. \textbf{This is the table the claim lives in.} \\[3pt]
\file{exp7\_conv\_order\_boxplot.pdf} &
Order per rule as a boxplot, with dashed lines at 1 and 2 so linear and quadratic
are readable at a glance. \\[3pt]
\file{exp7\_conv\_order\_scatter.pdf} &
Order against $k_\star$. If the claim holds this is a horizontal band: the vertical
position (rate) is common, the horizontal spread (when inactivity arrives) is not. \\
\bottomrule
\end{longtable}

\subsection{What the archived run actually shows}

From \file{exp\_2026-08-19\_16-28-37\_convergence\_rate}:

\begin{lstlisting}
rule               fitted    short    never   median order      median k*
RDelta                  3        0        0          1.462              2
RStep                   1        2        0          1.445              1
RDFO                    3        0        0          1.444              1
RGrad                   1        2        0          1.415              0
RGradCapped             1        2        0          1.415              0
RAdaptiveStep           1        2        0          1.445              1
RAdaptiveGrad           3        0        0          1.611              6
\end{lstlisting}

Read as intended: the median order sits in a narrow band, $1.415$ to $1.611$,
while median $k_\star$ ranges from $0$ to $6$. That is the predicted pattern --
common rate, different onset.

\finding{But the table must not be reported as it stands.} Three observations,
in increasing order of seriousness.

\begin{enumerate}[leftmargin=7mm,itemsep=3pt]
  \item \textbf{The problem set is three problems} -- AKIVA, BEALE, BOXBODLS --
        because the archived config records \jl{min\_var = 2} and \jl{max\_var = 2}.
        \file{config.jl} anticipates exactly this and forbids it in a comment on
        \jl{PROBLEM\_LIMIT}: \emph{``Three two-variable problems cannot resolve an
        asymptotic claim \dots\ Drop to 3 deliberately when debugging a script, and
        never report a table produced at that setting.''} The committed default is
        \jl{MAX\_VAR = 200}; the archive was produced with it narrowed to 2.
  \item \textbf{Four of the seven rules have a single fitted run.} A ``median'' over
        one observation is that observation. The \jl{short} column shows it plainly
        (2 of 3 runs discarded for four rules), which is the column doing its job,
        but it means the narrow band in column four is not yet evidence.
  \item \textbf{The archived config misstates the tolerance.} \jl{save\_config}
        records \jl{SOLVER\_PARAMS}, whose \jl{tol} is $10^{-5}$, but the script
        constructs and uses a local \jl{TRParams} with \jl{tol = 1e-10}. The
        archive therefore documents a run that did not happen. This matters more
        here than elsewhere, because the tolerance override is the one thing that
        makes the asymptotic measurement possible.
\end{enumerate}

\section{Experiment 9: first- versus second-order anchoring}

\subsection{The four claims}

The header states four, and the script is organised around them.

\begin{longtable}{@{}p{12mm}p{134mm}@{}}
\toprule
\endhead
\textbf{(a)} &
\textbf{The anchor collapse.} Started so the trajectory runs into a saddle, a
$\gnorm$-anchored rule reports the radius it would use at a solution. \jl{RGrad}
sets $\Delta = \mu\gnorm \to 0$ and halts. \jl{RDFO} finds $\Delta > \zeta\gnorm = 0$
at every iteration and contracts geometrically. Under
$\tau = \max\{\gnorm, -\lambda_{\min}\}$ both hold the radius at the scale of the
available curvature and escape. \\[3pt]
\textbf{(b)} &
\textbf{The measure gap.} The ratio $\tau_k/\gk$ is $1$ wherever the model is
convex and blows up near a saddle, which is exactly where a first-order diagnostic
reports success. \\[3pt]
\textbf{(c)} &
\textbf{The subsolver is half the story.} $\tau$ keeps the radius positive; it does
not make the step go anywhere. \\[3pt]
\textbf{(d)} &
\textbf{The model is the other half.} Over a positive semidefinite model
$\lambda_{\min}\ge 0$ identically, so $\tau \equiv \gnorm$ and the apparatus is an
expensive no-op that will still report \jl{:second\_order} at a saddle. \\
\bottomrule
\end{longtable}

\subsection{The test problems}

The primary problem is chosen so the failure needs no tuning to provoke:
\[
  f(x,y) \;=\; \tfrac{1}{4}x^4 - \tfrac{1}{2}x^2 + \tfrac{1}{2}y^2 ,
  \qquad x_0 = (0,\,0.7).
\]
There is a saddle at the origin with $\lambda_{\min} = -1$ and minimisers at
$(\pm 1, 0)$ with $f = -0.25$. The decisive property is that
$\partial f/\partial x = x^3 - x$ vanishes when $x = 0$, so the $x$-axis is
invariant under the gradient flow: a start on the $y$-axis runs \emph{exactly} into
the saddle. No parameter has to be tuned to make the pathology appear, which is
what makes it a fair test rather than a constructed one.

A second problem, the monkey saddle $x^3 - 3xy^2 + z^2/2$ from $(0.05, 0, 0.6)$,
guards against the finding being an artefact of one hand-built quartic. Its origin
is a degenerate critical point with a cubic escape direction.

The helper \jl{limit\_point(x)} classifies where a run stopped into
\jl{"saddle"}, \jl{"minimiser"} or \jl{"other"} by distance, which is what turns a
trajectory into a table entry.

\subsection{Parts (a) and (c): the measure $\times$ subsolver grid}

\jl{CONFIGS} is a $2\times2\times2$ grid: two rule families (\jl{RGrad},
\jl{RDFO}) $\times$ two measures ($\gnorm$, $\tau$) $\times$ two subsolvers
(\jl{SteihaugCG}, \jl{EigenPoint(SteihaugCG())}). Each entry also carries the
\jl{tol\_H} it should run with: $-1$ disables the second-order stopping test for
the $\gnorm$ rows, and $10^{-6}$ enables it for the $\tau$ rows. That coupling is
correct -- a second-order stopping test over a first-order measure would be
incoherent -- but it does mean measure and stopping test vary together, so the grid
isolates the \emph{pair}, not the measure alone.

\jl{anchor\_grid()} runs all eight through \jl{tr\_solve} directly rather than
through the harness. The file explains why: each is a statement about a single
trajectory on one hand-built problem, and routing it through the record layer would
add machinery without adding information.

\subsection{Part (b): the measure gap}

Three figures carry it, and the separation between them is deliberate:

\begin{itemize}[leftmargin=6mm,itemsep=2pt]
  \item \jl{plot\_radius} -- $\Delta_k$ per configuration on a log scale: the
        collapse under $\gnorm$, the plateau under $\tau$.
  \item \jl{plot\_measure\_gap} -- $\tau_k$ and $\gk$ on the same axes.
  \item \jl{plot\_gap\_ratio} -- the ratio $\tau_k/\gk$ on its own axis, with a
        dotted line at $1$. The docstring justifies the duplication: the ratio is
        the quantity with the interpretation, since where it exceeds $1$ the
        gradient alone calls the point critical and the curvature does not.
  \item \jl{plot\_curvature} -- $\lambda_{\min}(B_k)$, negative exactly where the
        first-order measure misleads.
\end{itemize}

\subsection{Part (d): the blind spot}

\jl{blind\_table()} holds the rule fixed at \jl{RGradTau} and the subsolver at
\jl{EigenPoint}, and varies only the model over \jl{ExactHessian}, \jl{SR1Model}
and \jl{LBFGSModel}. Because L-BFGS enforces $B \succ 0$ by construction,
$\lambda_{\min} > 0$ always, so $\tau \equiv \gnorm$ and the second-order variant
is identical to its first-order twin -- while still able to report
\jl{:second\_order}, because that status is a claim about $B$ and not about
$\nabla^2 f$. The file calls this the cheapest demonstration in the suite that a
status is only as strong as the model behind it.

\subsection{Part (e): the paired sweep}

\jl{paired\_sweep} runs every $\gnorm$-anchored rule immediately beside its
$\tau$-twin through the shared \jl{run\_experiment}, so these runs are archived and
resumable like every other experiment. The table to read is
\jl{success\_table(...; second\_order = true)}, whose last column counts
\emph{certified} \jl{:second\_order} stops. The reading rule is stated crisply: a
configuration solving every problem while certifying none has been running a
first-order method under a second-order name.

Two harness changes were needed for this part, and the file records both because
each would have inverted the conclusion:

\begin{itemize}[leftmargin=6mm,itemsep=2pt]
  \item \jl{solved} now counts \jl{:second\_order}. It previously tested
        \jl{=== :first\_order}, so a $\tau$-run that certified a minimiser scored as
        a \emph{failure}. The columns doing best would have shown zero reliability.
  \item \jl{RunRecord} now carries \jl{tau\_traj} and \jl{lambda\_traj}. They were
        dropped, so claim (b) could not be plotted from an archived run at all.
\end{itemize}

\subsection{What the archived run actually shows}

From \file{exp\_2026-08-19\_16-28-45\_second\_order}:

\begin{lstlisting}
configuration                      status   iters       limit        f(x*)      Delta_end
RGrad  |g|  Steihaug          first_order       1      saddle     0.000000     0.000e+00
RGrad  |g|  EigenPoint        first_order       1      saddle     0.000000     0.000e+00
RGrad  tau  Steihaug              stalled       2      saddle     0.000000     2.000e+00
RGrad  tau  EigenPoint       second_order       3   minimiser    -0.250000     0.000e+00
RDFO   |g|  Steihaug          first_order       1      saddle     0.000000     5.000e-01
RDFO   |g|  EigenPoint        first_order       3   minimiser    -0.250000     5.000e-01
RDFO   tau  Steihaug              stalled       2      saddle     0.000000     2.000e+00
RDFO   tau  EigenPoint       second_order       2   minimiser    -0.250000     1.000e+00
\end{lstlisting}

Claim (a) is confirmed sharply: both $\gnorm$-anchored \jl{RGrad} rows stop after
\emph{one} iteration at the saddle with $\Delta_{\text{end}} = 0$, reporting
\jl{:first\_order}. The table's own footnote makes the right point -- the run has
not failed by its own test, since $\gnorm$ really is zero there. It is the test
that is wrong.

Both $\tau$ + Steihaug rows end \jl{:stalled} at the saddle with the radius held at
$2.0$: $\tau$ kept the radius alive, and Steihaug had no reliable negative-curvature
step to spend it on. That is claim (c)'s first half, cleanly.

\finding{Claim (c) is contradicted by one row.} The header says only $\tau$
\emph{and} \jl{EigenPoint} together reach a minimiser. But
\jl{RDFO \ \textbar g\textbar \ EigenPoint} reaches the minimiser at $f = -0.25$
with the first-order measure. For \jl{RDFO}, \jl{EigenPoint} alone was sufficient
on this problem. The claim holds for \jl{RGrad} -- whose radius collapses to
exactly zero, leaving nothing for any subsolver to work with -- but not as a
general statement about both families. Either the claim should be narrowed to
\jl{RGrad}, or the difference between the two families should be explained.

\subsection{The blind-spot table is sharper than advertised}

\begin{lstlisting}
model                      status   iters       limit   min lambda_min   certified?
ExactHessian         second_order       3   minimiser         -1.0000          yes
SR1Model             second_order       1      saddle          1.0000          yes
LBFGSModel           second_order       1      saddle          1.0000          yes
\end{lstlisting}

The intended finding is the \jl{LBFGSModel} row: a certified \jl{:second\_order}
status at the saddle, from a model structurally incapable of seeing negative
curvature.

\finding{But \jl{SR1Model} behaves identically here}, and \file{config.jl} lists
SR1 among the models that \emph{can} report $\lambda_{\min} < 0$. The reason is
visible in the iteration count: the run stops after one iteration, and an SR1 model
initialised at the identity has not yet absorbed a single secant pair, so it
reports $\lambda_{\min} = 1$. The blind spot is therefore not only structural but
also transient -- a model that could see the curvature has not yet had the chance.
That is arguably a stronger result than the one claimed, and it is currently
unremarked.

\subsection{The paired sweep is thin}

\begin{lstlisting}
configuration              solved     median       mean  tail active  2nd order
RGrad                      1/2          20.0       20.0        1.000     1/2
RGradTau                   1/2           3.0        3.0        1.000     1/2
RGradCapped                0/2            --         --           --     0/2
RGradCappedTau             1/2           3.0        3.0        1.000     1/2
RDFO                       1/2           2.0        2.0        0.500     1/2
RDFOTau                    1/2           2.0        2.0        0.500     1/2
RAdaptiveGrad              0/2            --         --           --     0/2
RAdaptiveGradTau           1/2           2.0        2.0        0.500     1/2
\end{lstlisting}

With only two problems every cell is $0/2$, $1/2$ or $2/2$, and no configuration
exceeds $1/2$. The within-pair comparison is still legible -- \jl{RGrad} takes 20
iterations against \jl{RGradTau}'s 3, and \jl{RGradCapped} solves nothing while its
$\tau$-twin solves one -- but the counts cannot support a reliability statement.
More saddle problems are needed before this table carries weight in the manuscript.

\section{Shared machinery both files rely on}

\begin{longtable}{@{}p{40mm}p{106mm}@{}}
\toprule
\textbf{component} & \textbf{role} \\
\midrule
\endhead
\jl{RunRecord} &
One archived run: status, counts, and every trajectory, including
\jl{tau\_traj} and \jl{lambda\_traj} for second-order runs. \\[3pt]
\jl{RecordView} &
Dresses a \jl{RunRecord} as a solver result so the package's diagnostics
(\jl{inactivity\_index}, \jl{active\_fraction}, \jl{observed\_order}, \dots) apply
unchanged to an archived run. Its docstring gives the reason: without it every
experiment re-implements the same tail arithmetic, which is how a run and a table
come to disagree. \\[3pt]
\jl{run\_experiment} &
Runs the (problem $\times$ config) grid, archiving and skipping work already
present so a campaign can resume. \\[3pt]
\jl{solved} &
\jl{:first\_order} \emph{or} \jl{:second\_order}. \\[3pt]
\jl{active\_fraction} &
Share of the last \jl{tail} of a run on which the constraint bound. Used by
experiment 5 through \jl{RecordView}. Its docstring records the prediction that no
intermediate value survives in the limit, so one is a finding. \\[3pt]
\jl{inactivity\_index} &
$k_\star$, the first permanently inactive iteration, or \jl{nothing}. Used by
experiments 5 and 7. \\[3pt]
\jl{radius\_sums} &
$\sum\Delta_k$, $\sum\Delta_k^2$ and $\sum\Delta_k^2/M_k$, the last being \jl{NaN}
unless the run was traced with \jl{hessian\_norm = true}. \\[3pt]
\jl{ExperimentArchive} &
Timestamped directory with \file{experiment\_config.toml}, \file{tables/},
\file{figures/} and \file{data/}, closed by \jl{finalize\_archive}. \\
\bottomrule
\end{longtable}

\subsection*{A note on the two order estimators}

There are two, and they are not interchangeable:

\begin{itemize}[leftmargin=6mm,itemsep=2pt]
  \item \jl{estimate\_convergence\_order(grad\_traj; tail)} in
        \file{benchmark/harness.jl} takes a \emph{bare vector} and assumes the
        caller has already trimmed it. Its docstring warns about this in a boxed
        note. Experiment~7 uses this one, and does the trimming itself.
  \item \jl{observed\_order(stats; after\_inactivity, last, use)} in
        \file{src/Diagnostics/diagnostics.jl} takes a full result and does the
        trimming internally, and additionally skips rejected steps.
\end{itemize}

Experiment 7's hand-rolled \jl{inactive\_tail} predates \jl{RecordView} being
available to it, and duplicates logic \jl{observed\_order} already has. The two
also differ in a detail that could move a reported number: \jl{observed\_order}
filters out rejected steps before fitting, and \jl{estimate\_convergence\_order}
does not. Consolidating on the former, called through a \jl{RecordView}, would
remove the duplication and the discrepancy at once. Experiment 5 already takes
that route for \jl{active\_fraction} and \jl{radius\_sums}.

\section{Findings, collected}

\begin{longtable}{@{}p{8mm}p{60mm}p{74mm}@{}}
\toprule
& \textbf{observation} & \textbf{consequence} \\
\midrule
\endhead
1 & \textbf{exp5 and exp7} both ran on 3 two-variable problems &
\file{config.jl} explicitly forbids reporting a table at that setting. Re-run at
\jl{MAX\_VAR = 200} before either set of numbers is used. This is also the root
cause of finding 8. \\[3pt]
2 & exp7 records \jl{tol = 1e-5} in the archive but runs at \jl{1e-10} &
The archive documents a run that did not happen. \jl{save\_config} should receive
the local \jl{params}, not \jl{SOLVER\_PARAMS}. \\[3pt]
3 & exp9 claim (c) contradicted by the \jl{RDFO \textbar g\textbar\ EigenPoint} row &
That row reaches the minimiser without $\tau$. Narrow the claim to \jl{RGrad} or
explain the difference between the families. \\[3pt]
4 & exp9's SR1 row is blind for a different reason than L-BFGS &
Structural blindness versus a model that has not yet learned. Worth stating; it
strengthens the section. \\[3pt]
5 & exp9's paired sweep has two problems &
Every cell is $0/2$ or $1/2$. Add saddle problems before quoting reliability. \\[3pt]
6 & All three files end with a \jl{PROGRAM\_FILE} guard &
None can run standalone: \jl{RULES}, \jl{SOLVER\_PARAMS} and the harness come
from \file{initialisation.jl}. The guard, and the usage line in each header, promise
an entry point that does not work. This is the same defect that kept experiment 8
out of the suite. \\[3pt]
7 & exp7 duplicates \jl{observed\_order}'s trimming &
And differs from it by not skipping rejected steps. \\[3pt]
8 & exp5's \jl{measured} column hides a selection effect &
Run length is anti-correlated with the property measured, so four of six rules are
scored on a single, atypically slow run. The bias is conservative for the headline
claim but not for the intermediate values (\jl{RStep} $0.800$, \jl{RGrad} $0.400$).
Report $n$ beside each mean, or enlarge the problem set. \\[3pt]
9 & exp5 is the only script populating $\sum\Delta_k^2/M_k$ &
It passes \jl{solver\_kwargs = (hessian\_norm = true,)}. This is the column
\jl{tab: sums} in the manuscript leaves provisional on every row, so the machinery
to fill it already exists and is already exercised here. \\
\bottomrule
\end{longtable}

\vspace{4mm}
\noindent\textcolor{Muted}{\small Findings 1, 2 and 6 are defects in the
instrumentation rather than in the science. Findings 3 and 4 concern what the
manuscript may claim from experiment 9, and 5 and 8 concern how much weight the
corresponding tables can carry. Finding 9 is an opportunity rather than a defect.}

\end{document}
